\documentclass{article}

\usepackage{amsmath,amssymb}
\usepackage{xurl}
\usepackage[hidelinks]{hyperref}
\setlength{\emergencystretch}{2em}

\input{command}

\title{The False Converse in Wolf's Pauli-Block Truncation Proposition}
\author{}
\date{2026-08-23}

\begin{document}
\maketitle
\gapnote{false-source}{resolved}

\section{The printed equivalence}

In Section~2.4, after the Bloch-ball affine representation
$x\mapsto v+\Delta x$ in Eq.~(2.39), Wolf considers a Hermiticity-preserving
linear map $T:\MN{2}\to\MN{2}$. In the Pauli basis, write
\begin{align}
  \widehat T
  &=\begin{pmatrix}
      \widehat T_{00}&r^{\mathsf T}\\
      v&\Delta
    \end{pmatrix},
  &
  \widehat T_{ij}
  &=\frac12\tr\!\left(\sigma_iT(\sigma_j)\right).
\end{align}
Here $\Delta=(\widehat T_{ij})_{i,j=1}^{3}$ is the lower-right Pauli block,
and the row $r^{\mathsf T}$ and column $v$ are the two off-diagonal blocks.
Proposition~2.10 prints the claim that
\begin{align}
  \widehat T'
  :=\widehat T_{00}\oplus\Delta
  =\begin{pmatrix}
      \widehat T_{00}&0\\
      0&\Delta
    \end{pmatrix}
\end{align}
represents a positive, respectively completely positive, map if and only if
$\widehat T$ does.
\footnote{\cite[Section~2.4, Proposition~2.10, Eq.~(2.39)]{Wolf2012Quantum};
local source \path{Notes/WolfNoteTexSource/ch02_representations.tex},
lines~984--990.}

\section{What the proof establishes}

Let
\begin{align}
  \Theta(X)&=\tr(X)\Id-X,
  &D&=\mathrm{diag}(1,-1,-1,-1).
\end{align}
The Pauli transfer matrix of $\Theta$ is $D$, and
\begin{align}
  T'&=\frac{T+\Theta\circ T\circ\Theta}{2},
  &
  \widehat T'
  &=\frac{\widehat T+D\widehat T D}{2}
    =\widehat T_{00}\oplus\Delta.
  \label{eq:truncation}
\end{align}
Wolf's proof shows the two forward implications
\begin{align}
  T\text{ positive}
  &\Longrightarrow T'\text{ positive},
  \label{eq:positive-forward}\\
  T\text{ completely positive}
  &\Longrightarrow T'\text{ completely positive}.
  \label{eq:cp-forward}
\end{align}
It does not prove either converse. The proof route is exactly the averaging
in \eqref{eq:truncation}: double time reversal preserves positivity and
complete positivity in this setting, and these cones are convex.
\footnote{\cite[Proposition~2.10]{Wolf2012Quantum}; local source
\path{Notes/WolfNoteTexSource/ch02_representations.tex}, lines~992--998.}

The formalization follows this proof and states only the forward directions.
Pauli time reversal, its diagonal transfer coefficients, and the truncation are
\leanid{Wolf.pauliTimeReversal},
\leanid{Wolf.pauliTimeReversalDiagonal}, and
\leanid{Wolf.pauliBlockDiagonalTruncation}. The transfer-matrix identities are
\leanid{Wolf.pauliTransferEntry_timeReversal_comp} and
\leanid{Wolf.pauliTransferEntry_pauliBlockDiagonalTruncation}. The two
preservation theorems are
\leanid{Wolf.pauliBlockDiagonalTruncation_isPositiveMap} and
\leanid{Wolf.pauliBlockDiagonalTruncation_isCPMap}.
\footnote{All declarations are in
\doclink{QICLean/Channel/LorentzNormalForm/PauliBlockTruncation}.}

\section{Counterexample to both converses}

Define
\begin{align}
  T(X)&=\tr(X)\begin{pmatrix}3/2&0\\0&-1/2\end{pmatrix}.
  \label{eq:counterexample}
\end{align}
The map is Hermiticity preserving. It is not positive, because
\begin{align}
  T(\Id/2)
  &=\begin{pmatrix}3/2&0\\0&-1/2\end{pmatrix}
  \not\geq0.
\end{align}
Its Pauli transfer matrix has $\widehat T_{00}=1$, zero lower-right block
$\Delta=0$, and a nonzero entry only in the off-diagonal column $v$.
Consequently its truncation is
\begin{align}
  \widehat T'
  &=\mathrm{diag}(1,0,0,0),
  &
  T'(X)&=\frac{\tr(X)}2\Id.
\end{align}
This is the completely depolarizing channel and is completely positive.
Therefore $T'$ can be completely positive while $T$ is not even positive.
The reverse implication fails for positivity and, a fortiori, the reverse
implication fails for complete positivity.

\section{Conclusion}

The ``if and only if'' in Wolf's Proposition~2.10 is false as printed. The
proof in the source establishes the valid forward implications
\eqref{eq:positive-forward} and \eqref{eq:cp-forward}, and those are the
statements formalized in Lean. The explicit map \eqref{eq:counterexample}
shows that no converse can be inferred from Pauli-block truncation. The source
statement should therefore replace both equivalences by forward implications.

\bibliographystyle{alpha}
\bibliography{references}

\end{document}